\section{Related work}

\paragraph{Discrete choice models.} The conditional logit \citep{mcfadden1974} remains the
reference model, with the mixed logit \citep{mcfadden2000mixed} relaxing its independence
assumption. Recent work replaces parts of the utility with learned functions while keeping the
random-utility form: taste parameters as functions of covariates \citep{han2022tastenet},
alternative-specific architectures \citep{wang2020deep}, a learned term added to a specified
utility \citep{sifringer2020enhancing}, and one boosted ensemble per alternative under a shared
softmax \citep{salvade2024rumboost}. \citet{hillel2021systematic} benchmark these against
unconstrained classifiers and document how sensitive the ranking is to the split, which is the
sensitivity our protocol section addresses.

\paragraph{Language models as feature generators.} A separate line uses language models either as
direct predictors of behaviour \citep{brand2023using,horton2023large} or as generators of text
that a downstream model reads; the pipeline we study belongs to the second family. Our concern is
not whether such features help but what the reported help consists of, which places this work
closer to the literature on controls and leakage \citep{kapoor2023leakage} than to the modelling
literature.

\paragraph{Combining and calibrating.} Mixing predictive distributions is standard
\citep{wolpert1992stacked}, and stacking on out-of-fold predictions has known guarantees
\citep{vanderlaan2007super,yao2018using}. A linear pool is usually no worse than its components in
expectation, which is exactly why a mixture gain is weak evidence about the added component -- the
observation our uniform-floor control formalises. Temperature scaling
\citep{guo2017calibration} is the correction we apply to the conventional channel, fitted with the
same budget as the mixture weight so the baseline is not handicapped.
